\section{Introduction}
\label{sec:introduction}

When a graph intervention removes a circuit motif, what exactly has been
removed from the computation? Figure~\ref{fig:overview} illustrates why the
answer cannot be inferred from the edited structure alone. Two reconstructed
interventions on a fly connectome satisfy identical degree and module-pair
constraints and eliminate all common-relay triples for the designated sources
and sinks. Nevertheless, one retains both sources through separate relays,
while the other loses access to one. Their corresponding archived classifiers
achieve 97.45\% and 50.55\% accuracy. The structural certificates agree, but
the task information reaching the readout does not.

Degree-preserving rewiring and motif deletion isolate structural properties,
but a compensating edit can introduce replacement routes. The absence of a
motif therefore does not establish the absence of the associated computation.
The relevant object is the signal that the executed operator delivers to the
readout.

We address this gap with a task-information audit that separates three
questions: whether the declared structural constraints hold, whether the
information required by the task reaches the readout, and how effectively a
finite model uses that information. For the fixed two-step linear operator
studied here, propagating basis signals from the designated sources produces
a compact source-to-readout transfer. Its columns capture source access, and
their geometry describes how sources are mixed. Unlike a motif count, the
transfer includes every contributing route under the operator used by the
classifier.

The audit yields an exact result for the noiseless subspace XOR task: perfect
Bayes accuracy is possible if and only if both source columns are nonzero.
Neither a shared relay nor full transfer rank is required. Separate routes
can therefore preserve the task after a targeted common relay disappears.
Applied to the reconstructed MaleCNS interventions, the criterion
distinguishes the information-preserving and information-losing zero-relay
graphs, consistent with their archived outcomes.

A second distinction remains even when ideal task information survives.
Across two $k{=}8$ extractions, every nonzero dose retains both sources and
therefore preserves perfect Bayes solvability, yet finite-readout accuracy
rises from roughly 67\% to 94\%. Structural preservation, information
availability, and learned performance are different levels of evidence. On
visual parity and four-image CIFAR-10 probes, source mixtures and private
transfer directions further show why a visible source need not be equally
usable by a trained readout.

Our contributions are threefold:
\begin{itemize}
  \item We formulate a task-information audit that separates structural
        constraints, source-to-readout information, and finite-model
        behavior, with an exact criterion for a specified noiseless task.
  \item We identify a matched-connectome counterexample in which
        eliminating common relays can either preserve or remove the source
        information required by the task.
  \item We distinguish ideal information preservation from finite-readout
        performance and extend the audit to source recovery and mixing
        geometry in visual probes.
\end{itemize}
